%% Last Update: 17 Dec 2022
%

%% =================================================================================
\makeatletter
% تعریف یکسری متغیرها که کاربر می‌تواند بعدا آن ها را مقداردهی کند. 

\gdef\@type{نوع پروژه}
\def\type#1{\gdef\@type{#1}}

% عنوان محصول را تعیین می‌کند. این عنواند در ایجاد عنوان در مستند استفاده می‌شود این عنوان در هر مستند باید ایجاد شود در غیر این صورت از عنوان پیشفرض استفاده خواهد شد.
\gdef\@title{عنوان پروژه}
\def\title#1{\gdef\@title{#1}}

\gdef\@titleDesc{}
\def\titleDesc#1{\gdef\@titleDesc{#1}}

% زیر عنوان یک متن ساده را تعیین می‌کند که یک هدف مهم محصول را تعیین می‌کند این عنوان می تواند برای یک محصول در نظر گرفته نشود. از این داده برای  ایجاد عنوان و سایر مکان های محصول استفاده می‌شود.
\gdef\@subtitle{کاربرد محصول برای استفاده در زیر عنوان}
\def\subtitle#1{\gdef\@subtitle{#1}}
% افراد و گروه های که در تهیه این مستند و محصول همکاری داشته اند را تعیین می کند این داده همواره باید بیان شود. این داده در نوشتن عنوان و دیگر قسمت های مستند مورد استفاده قرار می‌گیرد.
\gdef\@author{}
\def\author#1{\gdef\@author{#1}}
\newcommand{\authorText}{\@author\,}

\gdef\@authorCode{}
\def\authorCode#1{\gdef\@authorCode{#1}}

% تاریخ نهایی نوشتن مستند را تعیین می‌کند این تاریخ در نوشتن عنوان استفاده می‌شود این تارخ باید تعیین شود در غیر این صورت به صورت پیش فرض یک تاریخ برای آن استفاده می شود.
\gdef\@date{ آخرین ویرایش: \myData}
\def\date#1{\gdef\@date{#1}}

\gdef\@supervisor{} 
\def\supervisor#1{\gdef\@supervisor{#1}}
\newcommand{\supervisorText}{\@supervisor\,}

\gdef\@adviser{نام استاد مشاور}
\def\adviser#1{\gdef\@adviser{#1}}

\gdef\@session{جلسه ارایه}
\def\session#1{\gdef\@session{#1}}

\gdef\@institute{پژوهشکده}
\def\institute#1{\gdef\@institute{#1}}

\gdef\@faculity{ نام دانشکده}
\def\faculity#1{\gdef\@faculity{#1}}

\gdef\@group{ نام گروه}
\def\group#1{\gdef\@group{#1}}

\gdef\@community{نام زیرگروه}
\def\community#1{\gdef\@community{#1}}

\gdef\@forwhat{پایان‌نامه یا رساله برای دریافت درجه}
\def\forwhat#1{\gdef\@forwhat{#1}}

\gdef\@field{در رشته مهندسی کامپیوتر گرایش}
\def\field#1{\gdef\@field{#1}}

\gdef\@version{اول}
\def\version#1{\gdef\@version{#1}}

% نام فایلی لوگوی مورد استفاده در نوشتار توسط این پارامتر مشخص می‌شود. 
\gdef\@logofile{\MyPath /images/logonotfound}
\def\logofile#1{\gdef\@logofile{#1}}
% اندازه فایل لوگوی موجود در متن توسط این پارامتر مشخص می‌شود.
\gdef\@logoScale{.3\textwidth}
\def\logoScale#1{\gdef\@logoScale{#1}}


%% ============================================================================================================================
%  برگه نخست در مستند را ایجاد می‌کند. در این برگه تنها نام پروردگار  ذکر می‌شود. برای ایجاد نام پرودگار یک تصویر استفاده می‌شود به نام god.ps  که باید در پوشه image کنار پرونده اصلی مستند قرار گرفته شده باشد.
\deflatinfont\besm[Scale=50]{Besmellah 2}
\DeclareDocumentCommand{\Godpage}{g g}{	
	\thispagestyle{empty}
	\begin{center}
	\IfValueTF{#1}{
		\IfValueTF{#2}{
			\includegraphics[width=#2\textwidth]{#1}
		}{%
			\includegraphics[width=.7\textwidth]{#1}
		}%
	}{%%
		\IfValueTF{#2}{
			{\besm #2}
		}{%
			{\besm b}

		}%
	}%%
	\end{center}
	\newpage
}%

\DeclareDocumentCommand{\Godpage}{g}{	
	\thispagestyle{empty}
	\begin{center}
		\IfValueTF{#1}{
			{\besm #1}
		}{%
			{\besm f}

		}%
	\end{center}
	\newpage
}%

\DeclareDocumentCommand{\GodpageFont}{gg}{	
	\thispagestyle{empty}
	\begin{center}
		\IfValueTF{#1}{
			\IfValueTF{#2}{
			}{%
			}%
			{\besm بسم‌الله الرحمن الرحیم}
		}{%
			{\besm بسم‌الله الرحمن الرحیم}

		}%
	\end{center}
	\newpage
}%


%% ============================================================================================================================

\newcommand{\pejoheshStyle}{
% این یک استایل ساده و رسمی.
% توسط این دستور تمامی footer و header های صفحه را حذف می‌کنیم. 
	\thispagestyle{empty}
	\begin{center}
		\includegraphics[width=\@logoScale]{\@logofile}\\*[10pt]
		\logofontR{\@institute}\medskip
		\vspace*{\stretch{2}}
		\typefontR\textbf{\@type}\\*[35pt]
		\titlefontR\textbf{\@title}\medskip
		\vspace*{\stretch{2}}
		\datafontR{\@author}\medskip
		\vspace*{\stretch{2}}
		\datafontR{\@date}
		\vspace*{\stretch{1}}
	\end{center}
	\clearpage
} %


%% ============================================================================================================================
% توسط دستورات defpersianfont و deflatinfont به ترتیب می توان یکسری فونت فارسی و انگلیسی دیگر تعریف کرد که در جاهای دیگر متن بتوان از آن استفاده کرد. برای استفاده کافی است که عبارتی که می خواهیم فونت آن عوض شود را به صورت زیر به عنوان نمونه بنویسیم.
\defpersianfont\titlefont[Scale=2.4]{IRTitr}
\newcommand{\lshortStyle}{
% استایلی شبیه به استایل شروع کتاب "مقدمه‌ای نه چندان کوتاه بر Latex"
	\definecolor{authorcol}{rgb}{.51,0,.51}
	\begin{flushleft}
	\vspace*{\stretch{.1}}
	\begin{flushright}\includegraphics[width=\@logoScale]{\@logofile}\end{flushright}
	\newlength{\centeroffset}
	\setlength{\centeroffset}{-0.5\oddsidemargin}
	\addtolength{\centeroffset}{0.5\evensidemargin}
	\addtolength{\textwidth}{-\centeroffset}
% توسط این دستور تمامی footer و header های صفحه را حذف می‌کنیم. 
	\thispagestyle{empty}
	
	\vspace*{\stretch{2}}
	
	\noindent\hspace*{\centeroffset}\makebox[0pt][l]{
		\begin{minipage}{\textwidth}
			\flushleft
			{ \titlefont\textcolor{magenta}\@title \\*[10pt]}
			\noindent\color{gray}{\rule[-1ex]{\textwidth}{5pt}\\*[4mm]
			\hfill{\large \bfseries\@type }}
		\end{minipage}
	}%
	
	\vspace{\stretch{2}}
	
	\noindent\hspace*{\centeroffset}\makebox[0pt][l]{
		\begin{minipage}{\textwidth}
			{\flushleft\textcolor{authorcol}{\bfseries\@author\\*[5pt]}}
			{\flushleft\textcolor{authorcol}{\bfseries\@supervisor\\*[5pt]}}
			{\flushleft\textcolor{authorcol}{\bfseries\small\@date}\\}
		\end{minipage}
	}%
		
	\vspace{\stretch{1.5}}

	\end{flushleft}
	\clearpage
}%

%% ===============================================================================

%\newfontfamily\sols[Scale=2.5,Script=Parsi,Language=Parsi,Mapping=farsidigits]{IranNastaliq}
\newcommand{\presStyle}[1]{
	\defpersianfont\titr[Scale=2]{B Nazanin Outline}

% ایجاد عنوان برای فایل ارایه (Presentation) 
% توسط این دستور تمامی footer و header های صفحه را حذف می‌کنیم. 
	\thispagestyle{empty}
	\deflatinfont\salam[Scale=30]{XB Kayhan}
	\begin{titlepage}
		\centering
		{\salam \textcolor{Blue}{#1}}\\
	   \begin{center}\includegraphics[width=\@logoScale\linewidth]{\@logofile}\end{center}
		\vspace{\stretch{2}}
		\textcolor{Green}{\huge\titr{\@title}}\\
		\vspace{\stretch{2}}
		\textcolor{cyan}{\@session}\\*[14pt]
		\textcolor{cyan}{\@author}\\*[14pt]
		\textcolor{Orange}{\@date}
		\vspace{\stretch{1}}
	\end{titlepage}
	\clearpage\newpage
	\pagestyle{empty}
	\pagenumbering{arabic}
}%	
	

%% ==========================================================================
%استایل مربوط به پایان‌نامه با دو لوگو
\newcommand{\thesisStyleOO}[2]{
	\ignorespaces
	\begin{center}
	\thispagestyle{empty}
	\begin{table}
		\begin{tabular}{ccc}
			\includegraphics[width=.18\textwidth]{#1}
			&
			\begin{minipage}{0.55\linewidth}
				\vskip 0.6cm
				\begin{center}
					\@institute\\* [0.4cm]
					\@faculity\\ [0.4cm]
					\@group \\*[1cm]
				\end{center}
			\end{minipage}
			&
			\includegraphics[width=.2\textwidth]{#2}
		\end{tabular}
	\end{table}
	\vspace*{\stretch{1}}
	\textbf{\LARGE{\@title}}
	\vspace*{\stretch{1}}
	\vskip 1cm
	\large{\@forwhat}\\[.4cm]
	\large{\@field}
	\vspace*{\stretch{1}}
	\vskip 1cm
	\Large{\textbf{\@author}}
	\vskip 1.5cm
	\Large{
	استاد راهنما:}
	 \\ [0.4cm] \Large{\@supervisor}
	\vskip 2cm
	\large{\@date}
	\end{center}
	\clearpage
}%

%% ==============================================================================
%استایل مربوط به پایان‌نامه با یک لوگو
\newcommand{\thesisStyleO}[1]{
\ignorespaces
\begin{center}
\thispagestyle{empty}
\begin{table}
\begin{tabular}{ccc}
& \includegraphics[width=\@logoScale]{\@logofile} & \\
&
\begin{minipage}{0.55\linewidth}
\vskip 0.6cm
\begin{center}
%\@institute\\* [0.4cm]
\@faculity\\ [0.4cm]
\@group \\*[1cm]
\end{center}
\end{minipage}
&
\end{tabular}
\end{table}
\vspace*{\stretch{1}}
\textbf{{\fontsize{15pt}{50pt}\selectfont\@title}}
\vspace*{\stretch{1}}
\vskip 1cm
\large{\@forwhat}\\[.4cm]
	\large{\@field}
	\vspace*{\stretch{1}}
	\vskip 1cm
	\Large{\textbf{\@author}}
	\vskip 1.5cm
	\Large{
	استاد راهنما:}
	 \\ [0.1cm] \Large{\@supervisor}
	\vskip 2cm
	\large{\@date}
	\end{center}
	\clearpage
}%

%% ================================================================================
%استایل مربوط به تکالیف
\newcommand{\homeworkS}{
\thispagestyle{empty}
\begin{center}
\begin{flushright}
\includegraphics[width=\@logoScale]{\@logofile}
\end{flushright}
\vspace{3cm}
{\farsifontshafighT \bfseries {\@title} }
\begin{tikzpicture}
\draw[red,line width=0.15cm](0,0)--(12,0);
\end{tikzpicture} 
\vspace{2cm}
\begin{minipage}{0.8\textwidth}
\begin{flushright} \large
\emph{\textcolor{Plum}{\@author}} \\[0.6cm]
\end{flushright}
\end{minipage}\\[1.5cm]
\begin{minipage}{0.8\textwidth}
\begin{flushright} \large
\emph{\@supervisor}\\[0.6cm]
\end{flushright}
\end{minipage}
\vfill\large \today
\end{center}
\clearpage
}

%% ==================================================================================

\let\oldmaketitle\maketitle
\renewcommand{\maketitle}{
\ifthenelse{\equal{\@titleStyle}{pejohesh}}{\pejoheshStyle}{ %1
\ifthenelse{\equal{\@titleStyle}{lshort}}{\lshortStyle}{ %2
\ifthenelse{\equal{\@titleStyle}{proposal}}{\proposalStyle}{ %3
\ifthenelse{\equal{\@titleStyle}{presentation}}{\presStyle}{ %4
\ifthenelse{\equal{\@titleStyle}{thesis}}{\thesisStyleO}{ %5
\ifthenelse{\equal{\@titleStyle}{homework}}{\homeworkS}{ %6
\oldmaketitle
} %6
} %5
} %4
} %3
}%2
}%1
}%maketitle


\newcommand\titlepagedecoration{%
\renewcommand\epigraphflush{flushright}
\renewcommand\epigraphsize{\normalsize}
\setlength\epigraphwidth{0.7\textwidth}

\definecolor{titlepagecolor}{cmyk}{1,.60,0,.40}

% The following code is borrowed from: https://tex.stackexchange.com/a/86310/10898
\begin{tikzpicture}[remember picture,overlay,shorten >= -10pt]

\coordinate (aux1) at ([yshift=-15pt]current page.north east);
\coordinate (aux2) at ([yshift=-410pt]current page.north east);
\coordinate (aux3) at ([xshift=-4.5cm]current page.north east);
\coordinate (aux4) at ([yshift=-150pt]current page.north east);

\begin{scope}[titlepagecolor!40,line width=12pt,rounded corners=12pt]
\draw
  (aux1) -- coordinate (a)
  ++(225:5) --
  ++(-45:5.1) coordinate (b);
\draw[shorten <= -10pt]
  (aux3) --
  (a) --
  (aux1);
\draw[opacity=0.6,titlepagecolor,shorten <= -10pt]
  (b) --
  ++(225:2.2) --
  ++(-45:2.2);
\end{scope}
\draw[titlepagecolor,line width=8pt,rounded corners=8pt,shorten <= -10pt]
  (aux4) --
  ++(225:0.8) --
  ++(-45:0.8);
\begin{scope}[titlepagecolor!70,line width=6pt,rounded corners=8pt]
\draw[shorten <= -10pt]
  (aux2) --
  ++(225:3) coordinate[pos=0.45] (c) --
  ++(-45:3.1);
\draw
  (aux2) --
  (c) --
  ++(135:2.5) --
  ++(45:2.5) --
  ++(-45:2.5) coordinate[pos=0.3] (d);   
\draw 
  (d) -- +(45:1);
  
\node at(1,24)[]{\includegraphics[width=\@logoScale]{\@logofile}};
  
\end{scope}
\end{tikzpicture}%
}

\newcommand\ReportTitleStyleB{

\defpersianfont\titleFont[Scale=4.0]{B Kamran Outline}
\defpersianfont\subtitleFont[Scale=1.5]{B Tabassom}

\begin{titlepage}
\noindent
\vskip 40mm
\hskip 20mm {\titleFont%
\textcolor{Blue}{\@title}}
\vskip 8mm
\hskip 20mm 
\begin{minipage}{0.55\textwidth}
\setstretch{1.2}{\subtitleFont\noindent\@titleDesc}
\end{minipage}

\null\vfill
\vspace*{-1cm}
\noindent
\hfill
\vspace*{3cm}
\begin{minipage}{0.30\linewidth}
\begin{flushleft}
\@version\\*[1mm]\@group
\end{flushleft}
\end{minipage}
%
\begin{minipage}{0.02\linewidth}
\textcolor{Blue}{\rule{1.5pt}{38pt}}
\end{minipage}
\titlepagedecoration
\end{titlepage}
}

\makeatother
